% Enterprise LLM TechDoc LaTeX Template (ctexart)
\documentclass[UTF8,a4paper,11pt,fontset=fandol]{ctexart}

% === Core packages ===
\usepackage{geometry}
\usepackage{booktabs}
\usepackage{array}
\usepackage{tabularx}
\usepackage{graphicx}
\usepackage{amsmath,amssymb}
\usepackage{hyperref}
\usepackage{xcolor}
\usepackage{caption}
\usepackage{float}
\usepackage{enumitem}
\usepackage{fancyhdr}
\usepackage{tcolorbox}
\usepackage{tikz}
\usetikzlibrary{arrows.meta,positioning,shapes.misc,shapes.geometric}

\geometry{a4paper, top=2.3cm, bottom=2.3cm, left=2.2cm, right=2.2cm}
\setlength{\parskip}{0.35em}
\setlength{\parindent}{2em}
\renewcommand{\arraystretch}{1.25}

% === Colors ===
\definecolor{Primary}{HTML}{1A1A2E}
\definecolor{Accent}{HTML}{1F6FEB}
\definecolor{Muted}{HTML}{6B7280}
\definecolor{BoxBg}{HTML}{F8FAFC}

% === Hyperref ===
\hypersetup{
  colorlinks=true,
  linkcolor=Accent,
  urlcolor=Accent,
  citecolor=Accent,
  pdftitle={<<DOC_TITLE>>},
  pdfauthor={<<DOC_AUTHOR>>}
}

% === Header / Footer ===
\pagestyle{fancy}
\fancyhf{}
\lhead{\textcolor{Muted}{<<DOC_CUSTOMER>>}}
\rhead{\textcolor{Muted}{<<DOC_VERSION>>}}
\cfoot{\textcolor{Muted}{\thepage}}
\renewcommand{\headrulewidth}{0.3pt}
\renewcommand{\footrulewidth}{0pt}

% === Boxes ===
\newtcolorbox{infobox}[1]{
  colback=BoxBg,
  colframe=Accent!45,
  title=\textbf{#1},
  boxrule=0.6pt,
  arc=2pt,
  left=6pt,right=6pt,top=6pt,bottom=6pt
}

% === Helpers ===
\newcolumntype{L}[1]{>{\raggedright\arraybackslash}p{#1}}
\newcolumntype{Y}{>{\raggedright\arraybackslash}X}

\title{\textbf{<<DOC_TITLE>>}\\[6pt]\large\textcolor{Muted}{<<DOC_SUBTITLE>>}}
\author{<<DOC_AUTHOR>>}
\date{<<DOC_DATE>>}

\begin{document}
\maketitle

\begin{infobox}{文档控制信息}
\begin{tabularx}{\linewidth}{L{3.2cm}Y L{3.0cm}Y}
\toprule
客户名称 & <<DOC_CUSTOMER>> & 行业 & <<DOC_INDUSTRY>> \\
版本 & <<DOC_VERSION>> & 日期 & <<DOC_DATE>> \\
作者 & <<DOC_AUTHOR>> & 审阅 & <<DOC_REVIEWER>> \\
保密级别 & <<DOC_CONF>> &  &  \\
\bottomrule
\end{tabularx}
\end{infobox}

\vspace{0.5em}
\tableofcontents
\newpage

\section{概述}
\subsection{目标与范围}
本文档用于描述\textbf{目标大模型服务（以合同与实际接入平台为准）}能力与客户现有产品/系统的结合方式，输出可交付的技术方案结构（排版与内容为模板样例）。对任何具体能力、指标与 SLA 的承诺以合同与实际开通能力为准。

\subsection{术语与约束}
\begin{itemize}[leftmargin=2.2em]
  \item 大模型能力范围：仅覆盖已开通/已采购能力；不确定项标记“待确认”。
  \item 数据合规：遵循数据分级分类、最小权限、脱敏与审计要求。
  \item 输出边界：本模板强调“能力—系统—场景—KPI”的可追溯关系，避免泛化营销表述。
\end{itemize}

\section{客户现有产品与集成切入点}
\begin{infobox}{现有产品清单（示例）}
\begin{tabularx}{\linewidth}{L{3.5cm} L{2.7cm} Y}
\toprule
产品/系统 & 责任部门 & 现状与说明 \\
\midrule
<<CURRENT_PRODUCTS_TABLE>>
\bottomrule
\end{tabularx}
\end{infobox}

\section{大模型能力描述（按可交付能力组织）}
\begin{infobox}{能力清单（示例，具体以实际开通为准）}
\begin{tabularx}{\linewidth}{L{5.2cm} Y}
\toprule
能力 & 说明与典型用途 \\
\midrule
<<LLM_CAP_TABLE>>
\bottomrule
\end{tabularx}
\end{infobox}

\section{能力与现有产品结合分析（映射与收益）}
\subsection{映射方法}
建议采用“\textbf{场景}→\textbf{能力}→\textbf{集成点}→\textbf{收益}→\textbf{KPI}”的结构化映射，便于后续验收与运营复盘。

\subsection{能力—系统—场景映射表（模板）}
\begin{infobox}{映射表（示例）}
\begin{tabularx}{\linewidth}{L{3.3cm} L{3.2cm} L{3.4cm} Y}
\toprule
场景 & 大模型能力 & 集成系统/接口 & 预期收益与KPI \\
\midrule
<<MAPPING_TABLE>>
\bottomrule
\end{tabularx}
\end{infobox}

\section{业务场景方案与优势（按场景输出）}
<<SCENARIO_SECTIONS>>

\section{总体架构（示意）}
\begin{figure}[H]
\centering
\begin{tikzpicture}[
  node distance=12mm and 12mm,
  box/.style={draw=Primary!40, rounded corners=2pt, fill=BoxBg, align=left, inner sep=6pt},
  svc/.style={draw=Accent!60, rounded corners=2pt, fill=Accent!6, align=left, inner sep=6pt},
  arr/.style={-Latex, draw=Primary!55, line width=0.7pt}
]
  \node[box] (user) {用户/员工\\(APP/小程序/坐席)};
  \node[box, right=of user] (entry) {业务入口\\(客服工作台/IM/门户)};
  \node[svc, right=of entry] (orchestrator) {编排层\\(提示词/路由/工具调用)};
  \node[svc, below=of orchestrator] (llm) {大模型推理服务\\(对话/理解/生成)};
  \node[box, right=of orchestrator] (systems) {客户系统\\(工单/CRM/订单/知识库)};
  \node[box, below=of systems] (kb) {知识库与数据源\\(文档/FAQ/日志)};

  \draw[arr] (user) -- (entry);
  \draw[arr] (entry) -- (orchestrator);
  \draw[arr] (orchestrator) -- node[above,sloped]{调用} (systems);
  \draw[arr] (orchestrator) -- node[right]{推理/生成} (llm);
  \draw[arr] (kb) -- node[below,sloped]{检索增强} (orchestrator);
  \draw[arr] (systems) -- node[right]{事件/结果} (orchestrator);
\end{tikzpicture}
\caption{大模型能力与客户系统集成总体架构示意（模板）}
\label{fig:arch}
\end{figure}

\section{部署、安全与运维要点（模板）}
\subsection{部署模式与数据范围}
\begin{itemize}[leftmargin=2.2em]
  \item 部署模式：<<DEPLOY_MODE>>
  \item 主要数据源：<<DEPLOY_DATA_SOURCES>>
  \item 安全与合规：<<DEPLOY_SECURITY>>
\end{itemize}

\subsection{运维与运营指标}
\begin{itemize}[leftmargin=2.2em]
  \item SLA：<<OPS_SLA>>
  \item 监控项：<<OPS_MONITORING>>
  \item 成本模型：<<OPS_COST>>
\end{itemize}

\section{验收建议（模板）}
\begin{infobox}{建议的验收维度}
\begin{itemize}[leftmargin=2.2em]
  \item 功能验收：场景覆盖、工具调用链路、权限与审计、引用溯源（RAG）。
  \item 指标验收：自助解决率、时延、人工转接率、准确率（按业务定义）。
  \item 合规验收：敏感信息处理、越权防护、日志留存、数据出境（如适用）。
\end{itemize}
\end{infobox}

\end{document}
